\documentclass[11pt,letterpaper]{article}
\usepackage[T1]{fontenc}
\usepackage[utf8]{inputenc}
\usepackage{lmodern}
\usepackage[margin=1in,headheight=15pt]{geometry}
\usepackage{amsmath,amssymb,amsthm,mathtools,mathrsfs}
\usepackage{microtype,booktabs,longtable,array,enumitem}
\usepackage[dvipsnames]{xcolor}
\usepackage{fancyhdr,listings}
\usepackage[colorlinks=true,linkcolor=MidnightBlue,citecolor=MidnightBlue,urlcolor=MidnightBlue]{hyperref}
\usepackage[nameinlink,noabbrev]{cleveref}
\hypersetup{pdftitle={Exact Symbolic Computation with Surreal and Surcomplex Numbers},pdfauthor={Research and implementation study},pdfsubject={Hahn fields, computability, transseries, symbolic CAS, and Wolfram Language}}
\pagestyle{fancy}\fancyhf{}
\fancyhead[L]{\small Surreal and surcomplex computer algebra}
\fancyhead[R]{\small Representations and algorithms}
\fancyfoot[C]{\thepage}
\setlength{\parskip}{0.3em}\setlength{\parindent}{1.3em}
\setlist{nosep,leftmargin=2em}
\allowdisplaybreaks[2]\emergencystretch=2em
\newtheorem{theorem}{Theorem}[section]
\newtheorem{proposition}[theorem]{Proposition}
\newtheorem{lemma}[theorem]{Lemma}
\newtheorem{corollary}[theorem]{Corollary}
\theoremstyle{definition}\newtheorem{definition}[theorem]{Definition}
\newtheorem{example}[theorem]{Example}
\theoremstyle{remark}\newtheorem{remark}[theorem]{Remark}
\newcommand{\No}{\mathbf{No}}\newcommand{\SC}{\mathbf{SC}}
\newcommand{\R}{\mathbb R}\newcommand{\C}{\mathbb C}
\newcommand{\Q}{\mathbb Q}\newcommand{\N}{\mathbb N}\newcommand{\Z}{\mathbb Z}
\newcommand{\Acal}{\mathscr A}\newcommand{\HH}{\mathcal H}\newcommand{\OO}{\mathcal O}
\newcommand{\m}{\mathfrak m}\newcommand{\st}{\operatorname{st}}
\newcommand{\supp}{\operatorname{supp}}\newcommand{\lc}{\operatorname{lc}}
\newcommand{\Res}{\operatorname{Res}}\newcommand{\res}{\operatorname{res}}
\newcommand{\NF}{\operatorname{NF}}\newcommand{\Tr}{\operatorname{Tr}}
\newcommand{\id}{\operatorname{id}}\newcommand{\vH}{v_{\!H}}
\newcommand{\HInt}{\mathop{\int^{\!H}}\nolimits}
\newcommand{\code}[1]{\texttt{#1}}
\newcolumntype{L}[1]{>{\raggedright\arraybackslash}p{#1}}
\numberwithin{equation}{section}\setcounter{tocdepth}{2}
\lstdefinestyle{wl}{basicstyle=\ttfamily\small,breaklines=true,columns=fullflexible,
 keepspaces=true,showstringspaces=false,frame=single,rulecolor=\color{gray!40},
 backgroundcolor=\color{gray!3},xleftmargin=0.3em,xrightmargin=0.3em,
 aboveskip=0.7em,belowskip=0.6em}
\lstset{style=wl}
\begin{document}
\begin{titlepage}\centering
\vspace*{0.65cm}
{\Large\scshape A research and implementation study\par}
\vspace{0.7cm}
{\Huge\bfseries Exact Symbolic Computation\par}
\vspace{0.3cm}
{\LARGE with Surreal and Surcomplex Numbers\par}
\vspace{0.7cm}
{\Large Representable domains, support certificates,\par algorithms, and a Wolfram Language prototype\par}
\vspace{0.6cm}{\large September 21, 2026\par}
\vspace{0.6cm}
\begin{minipage}{0.94\textwidth}
\begin{abstract}
A useful computer algebra system for surreal numbers should not attempt to turn the entire proper class into one unrestricted numeric type. It should provide compatible, explicitly presented domains: finite Conway constructions; exact coefficient and exponent fields; sparse Hahn expressions; rational functions of ordered monomials; selected algebraic roots; certified lazy series; and suitable transseries and analytic-germ extensions. Surcomplex arithmetic then adds a quadratic extension, while surcomplex analysis also requires domains, branch choices, and an integration category.

We distinguish finite exact descriptions from canonical forms, computable coefficients, effective leading terms, decidable equality, and finite precision computations. A completely explicit rational-monomial domain already supports exact higher-rank arithmetic, comparison, standard parts, and conjugation. A finite-grid dependency theorem explains when strong Hahn formulae become coefficient algorithms; a halting reduction explains why unrestricted coefficient programs cannot carry a total zero-test. We translate the three supplied manuscripts into computational interfaces for local Taylor evaluation, radius-free Hahn germs, preparation, finite quotient algebras, residues, and contour periods. We also discuss Gonshor exponentiation, phase-dependent global surcomplex exponentials, current transseries zero-testing results, and why ordinary asymptotic expansions are not automatically exact surreal values.

An accompanying Wolfram Language package implements the rational-monomial core over $\Q(i)$ with lexicographic integer-vector valuations. The package and finite illustrative calculations are tested separately from the unimplemented general analytic and transseries layers.
\end{abstract}
\end{minipage}
\vfill
\begin{minipage}{0.94\textwidth}\small
\textbf{Status and provenance.} This is a design article with proofs of elementary computational claims and a deliberately limited prototype. Results drawn from the three supplied research manuscripts are explicitly attributed to those manuscripts, not promoted to independently verified published theorems. Published background, implementation proposals, and executed tests are kept distinct. No claim of a complete surreal CAS or a universally decidable symbolic language is made.
\end{minipage}
\end{titlepage}
\begingroup\small\setlength{\parskip}{0pt}\tableofcontents\endgroup\newpage

\section{The answer is a hierarchy, not a single data type}\label{sec:intro}
\subsection{What a successful system should promise}
The most practical starting point is an exact ordered field of rational expressions in finitely many non-Archimedean monomials, followed by algebraic and carefully specified functional extensions. Such a system can manipulate genuine infinite and infinitesimal numbers without replacing them by limits, very large floating-point numbers, or truncated polynomials. It can also retain a finite description of an infinite Hahn expansion. For example,
\begin{equation}\label{eq:firstexample}
 \frac{1}{1-\omega^{-1}}
 =\sum_{n\ge0}\omega^{-n}
\end{equation}
is an exact rational expression as well as a strongly summable normal form. Computing with its left side requires neither enumerating the right side nor performing a limiting numerical process.

There is no single largest practically meaningful answer to ``which surreal numbers are exactly representable?'' A language allowing rational operations names one class; adding selected algebraic roots names a larger one; adding named analytic functions or programs names still more. Each enlargement changes the available decision procedures. Even two representations of the \emph{same mathematical subset} can have very different computational properties. The distinction is demonstrated rigorously in \cref{sec:effectivity}.

For surcomplex numbers, algebraic arithmetic is not the main additional obstacle. If a real domain is an effective ordered field $F$, pairs $(a,b)$ implement $F[i]$, with
\[
 (a+ib)^{-1}=\frac{a-ib}{a^2+b^2}\qquad(a+ib\ne0).
\]
The difficulty increases when a function rather than a field operation is requested. ``Evaluate an exponential'', ``differentiate'', ``take a limit'', and ``integrate around a curve'' each need a precise semantic interpretation. The supplied manuscripts show why one cannot silently identify fine-local analyticity, Hahn-coherent holomorphy, and contour integration \cite{SA,AG,CT}.

\subsection{The supplied manuscripts as computational specifications}
We use the following short names throughout. \textbf{SA} is \emph{Surcomplex Analysis: Infinitesimal Calculus, Hahn-Coherent Holomorphy, and Contour Theory}. \textbf{AG} is \emph{Infinitesimal Analytic Geometry over the Surcomplex Numbers: Radius-Free Hahn Germs, a Nullstellensatz, and Conservation of Singular Zeros}. \textbf{CT} is \emph{Jordan Separation, Cauchy Theory, and Stokes Formulae on the Surcomplex Plane} \cite{SA,AG,CT}.

SA supplies distinct function categories, common-domain coefficientwise calculus, and examples of nonunique global exponential extensions. AG supplies a radius-free coefficient-germ algebra and support-controlled formulae for division, finite-dimensional normal forms, and residues. CT supplies explicit categories of chains and microscopic contour realizations. In particular, AG itself warns that finite mathematical dependencies are not necessarily effectively enumerable. Our question is what extra representations and algorithms make those formulae executable.

The article does not re-audit AG's Noetherian, Jacobson, or deformation proofs, nor CT's general residue-transformation argument. Their computational consequences are conditional on the stated manuscript results \emph{and} on effective input operations. For polynomial examples, we can independently check the displayed algebraic identities without implementing the general germ theory.

\subsection{An initial capability map}
\begin{center}\small
\begin{longtable}{@{}L{0.22\textwidth}L{0.34\textwidth}L{0.35\textwidth}@{}}
\toprule Representation & Exact objects available & Principal computational boundary\\\midrule\endhead
Hereditarily finite number cuts & Dyadic rational numbers & Not closed under division; finite recursive trees can grow rapidly.\\[0.4em]
Finite Hahn sums & Finite normal forms with supported coefficients and exponents & Addition and multiplication; reciprocals generally require a larger representation.\\[0.4em]
Rational monomial expressions & A non-Archimedean fraction field & Total arithmetic, equality, and real order for effective parents; not algebraically closed.\\[0.4em]
Selected algebraic extensions & Algebraic roots over an effective base & Root selection and exact coefficient procedures are part of the representation.\\[0.4em]
Certified lazy Hahn series & Infinite supports given by suitable finite programs & Equality and leading-term extraction require extra structure; arbitrary programs do not suffice.\\[0.4em]
Transseries and differential-algebraic descriptions & Exponential--logarithmic and selected differential-equation extensions & Use actual effective transseries algorithms; not every surreal expression belongs to their input class.\\[0.4em]
Analytic germs and Hahn families & Local functions, finite jets, coherent periods & Domain, support, branch, and ordinary coefficient algorithms must be supplied.\\
\bottomrule
\end{longtable}\end{center}
The rows are not a claim that every object in an abstract Hahn field or every analytic germ is finitely describable. They identify workable \emph{presentation families} inside the mathematical objects.

\section{Exactness, size, and the meaning of a symbolic name}\label{sec:exactness}
\subsection{Six different promises}
An \emph{exact denotation} means that a finite expression has a specified mathematical value, without rounding. A \emph{canonical representation} means that equal values have identical normalized data. A \emph{coefficient procedure} computes a requested coefficient of a series. A \emph{decision procedure} terminates on every valid input for equality or order. A \emph{precision procedure} produces a certified approximation in a specified topology or filtration. A \emph{proof-carrying representation} additionally supplies checkable evidence for the conditions making an expression meaningful.

These promises are independent enough that they should be separate capabilities. A selected algebraic root is exact without being a radical expression. A convergent or formal series specified by a recurrence can be exact without a general equality decision. A finite jet is exact in a quotient algebra but does not identify one exact function. A coefficient oracle can compute every requested coefficient while being unable to determine whether all coefficients vanish.

Consider the following three outputs:
\[
 A=\frac1{1-t},\qquad B=1+t+\cdots+t^9+O_v(t^{10}),\qquad
 C=\operatorname{StrongSum}(n\mapsto t^n).
\]
Here $A$ names one exact element. $B$ names a precision class unless it carries an additional exact-tail specification. $C$ names the same element as $A$ only after the system has certified strong summability and the intended monomial interpretation. Printing an infinite summation symbol, or leaving a Wolfram expression inactive, does not itself provide that certificate.

\subsection{Why all surreals cannot be stored}
Finite strings over a finite or countable alphabet form a countable set. A parameter-free symbolic language interpreted by ordinary finite programs can therefore name only countably many individual values. Already the real numbers are too numerous; the proper class $\No$ is a still larger target. The obstruction is not that a computer must store an enormous but finite exponent. Most values simply have no name in a chosen effective language.

An opaque coefficient oracle changes the interface, not this conclusion for ordinary finite programs. If arbitrary real or set-valued oracles are admitted as external parameters, the resulting object is exact \emph{relative to the oracle}; it is not a finite standalone computable encoding. Likewise, a symbol $c$ with assumptions describes a parameter or a family until a particular value is specified.

A practical implementation should thus use set-sized \emph{parent domains} and explicit embeddings between them. It may use a class-sized theory to justify those parents without ever allocating a class-sized list, ring, or support. This agrees with the workspace convention in all three manuscripts \cite{SA,AG,CT}.

\subsection{Surreal values, ordinal values, and expressions for infinity}
The ordinal $\omega$, its surreal image, and a CAS symbol for an infinite limit are not interchangeable. On ordinal inputs, surreal addition agrees with natural, rather than ordinary noncommutative ordinal, addition. In particular, ordinary ordinal arithmetic gives $1+\omega=\omega$, whereas surreal field arithmetic gives $1+\omega=\omega+1>\omega$. A system supporting both needs distinct constructors and operations \cite{Gonshor}.

Similarly, Wolfram Language's documented \code{Infinity} participates in extended-limit arithmetic; it is not a named element of a non-Archimedean ordered field \cite{WLInfinity}. The intended surreal identity is
\[
 \omega-\omega=0,\qquad \omega^{-1}>0,
\]
not an indeterminate subtraction and not a reciprocal equal to zero. In the prototype, \code{Infinity} appears only as the conventional metadata value $v(0)=+\infty$, never as a stored surreal field element.

\section{Representing the numbers}\label{sec:representations}
\subsection{Conway cuts and sign expansions}
The foundational cut construction is mathematically natural and useful for validation, educational examples, simplicity questions, and conversion of small finite cases. A legal number cut has every left option smaller than every right option. Its value is the \emph{simplest} number between the options, not an arbitrary point in an interval. For instance $\{0\mid4\}=1$, not the midpoint $2$. A cut constructor cannot therefore be implemented by choosing a convenient solution of inequalities \cite{Gonshor}.

Hereditarily finite legal number cuts, equivalently finite sign expansions, represent exactly the dyadic rationals. Memoization and directed acyclic graphs can avoid repeated subtrees, but do not change that domain. Division immediately leaves it: $1/3$ is not dyadic. Finite cuts whose options are already infinite objects form a different, more expressive language, and require algorithms for simplicity relative to those option presentations.

Compressed infinite sign strings and cut schemas can name $\omega$, real numbers with effective digit descriptions, and many more values. The compression language must say how a transfinite block is indexed, how its legality is certified, and which comparisons terminate. An ordinal notation appearing in the index is not a universal algorithm for all ordinal or surreal questions. Birthdays and simplicity comparisons should consequently be optional capabilities, not prerequisites for ordinary field arithmetic.

\subsection{Coefficient and exponent parents}
Start with an effective coefficient field $k$. The strongest simple defaults are $\Q$, $\Q(i)$, real algebraic numbers, and complex algebraic numbers with exact root data. Finite number fields offer an intermediate performance choice. Wolfram's \code{Root}, \code{AlgebraicNumber}, and associated algebraic-number operations are natural coefficient tools \cite{WLRoot,WLAlgebraic}.

A larger coefficient expression language containing $\pi$, $e$, logarithms, or special-function values can still denote exact constants. It must not inherit the algebraic backend's guarantee of deciding every zero and every sign. The required distinction is not ``algebraic is exact, transcendental is approximate''; both can be exact. It is that their supported algorithms and semantic relations differ.

For exponents, an especially useful parent is
\begin{equation}\label{eq:lexgroup}
 \Gamma_d=\Q^d_{\mathrm{lex}},
\end{equation}
with exact vector addition and lexicographic comparison. A concrete surreal embedding sends
\[
 (a_1,\ldots,a_d)\longmapsto
 a_1\omega^{d-1}+a_2\omega^{d-2}+\cdots+a_d.
\]
The displayed powers and addition here are surreal operations. More general finite rational spans of certified ordered surreal exponents are possible. Arbitrary named generators must not simply be assumed independent: hidden exponent relations would destroy uniqueness of monomials and invalidate equality tests.

The coefficients and the exponents may both have recursive structure. For example, a monomial with exponent $\omega$ is finitely describable even though it is not in a rank-one rational-exponent lattice. A recursive exponent backend should export its own equality, order, and normalization procedures, with a well-founded construction rank. An arbitrary reference back to the same object is not a legitimate normal form.

\subsection{Hahn normal forms and their interpretation}
Fix an ordered set-sized subgroup $\Gamma\subset\No$. In the convention of the supplied manuscripts,
\begin{equation}\label{eq:hahn}
 k((t^\Gamma))=\left\{\sum_{\gamma\in S}c_\gamma t^\gamma:
 S\subset\Gamma\text{ well ordered}\right\},\qquad
 t^\gamma=\omega^{-\gamma}.
\end{equation}
The monomial notation uses Conway's omega-map. It should be internally distinguished from a general exponential definition of powers. In particular, an implementation must not introduce the rewrite $\omega^a=\exp_{\No}(a\log_{\No}\omega)$ for arbitrary surreal $a$ merely by borrowing the usual real-power notation. The two constructions need an actual compatibility theorem on the chosen argument class \cite{Gonshor,SA}.

For real coefficients, a nonzero series is positive exactly when its coefficient at the least exponent is positive. For complex coefficients there is a valuation and coefficientwise conjugation, but no field ordering. Its real and imaginary parts are obtained coefficientwise. When $\Gamma$ is divisible, the real and complex Hahn workspaces in the supplied manuscripts are respectively real closed and algebraically closed. The relevant Hahn-field algebraic-closure result is an existence theorem, not a promise about arbitrary support programs \cite{Poonen}.

A finite support can be stored as a sparse sorted map from normalized exponents to nonzero coefficients. Addition merges equal keys. Multiplication forms pairwise exponent sums and combines coefficients. If the coefficient zero-test is partial, the sparse map may contain unresolved terms and cannot honestly advertise its first stored key as the valuation.

\subsection{Finite sums versus rational monomial fields}
Finite Hahn sums form a group algebra, not a field. Their nonzero monomials are invertible, but the reciprocal of $1-t$ has infinite support. Instead of expanding every inverse, retain rational expressions in independent monomials. This yields a useful exact field with a finite symbolic representation for every element.

\begin{theorem}[Effective rational-monomial core]\label{thm:rationalcore}
Let $k$ be an effective field, with effective order when real comparisons are requested. Let $X_1,\ldots,X_d$ denote monomials with independent standard-basis exponents in $\Z^d_{\mathrm{lex}}$. Then $k(X_1,\ldots,X_d)$ embeds in $k((t^{\Z^d_{\mathrm{lex}}}))$. Rational expressions admit terminating field arithmetic and equality algorithms. The valuation and initial coefficient of a nonzero $P/Q$ are
\begin{equation}\label{eq:rationalvaluation}
 v(P/Q)=\min\supp P-\min\supp Q,\qquad
 \lc(P/Q)=\frac{\lc(P)}{\lc(Q)}.
\end{equation}
For a real input, its sign is the sign of this initial coefficient.
\end{theorem}
\begin{proof}
Distinct polynomial monomials have distinct exponent vectors, so a nonzero polynomial cannot vanish in the Hahn field. Write a nonzero polynomial denominator as $Q=c t^\beta(1+\epsilon)$, where the finitely many exponents of $\epsilon$ are strictly positive. The strong geometric series $\sum_{n\ge0}(-\epsilon)^n$ is a Hahn inverse; the positive-support lemma justifies each coefficient. Thus the polynomial embedding extends to the fraction field.

Finite polynomial arithmetic over $k$ computes sums, products, and reciprocals of fractions. Equality reduces to testing the finitely many coefficients of $PQ'-P'Q$. The leading product term cannot cancel because the least exponent in each factor is unique, proving \eqref{eq:rationalvaluation}. The Hahn ordering gives the final assertion. No infinite expansion is required by any of these algorithms.
\end{proof}
Polynomial gcd reduction, when supplied by the coefficient backend, yields reduced fractions; normalizing the denominator's leading coefficient to one removes scalar ambiguity. The prototype uses Wolfram's rational polynomial operations for $k=\Q(i)$. The theorem itself does not require gcd computation for equality.

Rational exponent vectors can be supported by refining a common denominator lattice before a finite computation. This permits $t^{1/2}$ and $t^{1/3}$ without confusing them with independent variables. A finite lattice refinement still does not make every algebraic root a finite monomial expression.

\subsection{Algebraic roots and Puiseux representations}
An algebraic extension can store a polynomial over its current parent together with data selecting one root: a real isolating interval or Thom encoding in an ordered extension, or a residual root and sufficient valuation separation in a valued presentation. A polynomial alone does not select one of its roots. Nor does an arbitrary integer root index inherit the meaning of Wolfram's ordinary complex root ordering at non-Archimedean coefficients \cite{WLRoot}.

One must not assume the base ordered field is dense in its real closure. If $T$ is positive infinite, no element of $\Q(T)$ lies between $\sqrt T$ and $\sqrt T+1$: a nonzero rational function of $T$ has an integer leading power, whereas every point of that interval has leading power $T^{1/2}$. Thus root separation may require Thom data or endpoints in a suitable extension, not intervals with endpoints confined to the original rational parent.

For a simple binomial root, the construction is explicit. If $a=c t^\gamma(1+\epsilon)$ with $v(\epsilon)>0$, then a chosen $n$th root is
\begin{equation}\label{eq:binomialroot}
 c^{1/n}t^{\gamma/n}
 \sum_{j\ge0}\binom{1/n}{j}\epsilon^j.
\end{equation}
It requires a coefficient-root choice and the exponent $\gamma/n$. The series is strongly summable. Multiplication by the $n$th roots of unity supplies the other roots. Formula \eqref{eq:binomialroot} is an exact lazy description, not usually a finite normal form.

In rank one and characteristic zero, Puiseux series are the classical algebraic enlargement of Laurent series over an algebraically closed coefficient field: the union of the fields $k((t^{1/n}))$ is algebraically closed \cite[Corollary 7]{Poonen}. For polynomial equations over an effective rational-function base, Newton--Puiseux expansion with residual factorization and branch selection is an appropriate implementation strategy. The full abstract Puiseux field still contains arbitrary, noncomputable coefficient sequences; its algebraic closedness does not encode those sequences.

Higher-rank algebraic algorithms should use nested valued fields, valuation extensions, and explicit selected-root objects, rather than flattening every scale into one real exponent. An effective real closure of the rational-monomial ordered field, followed by adjoining $i$, is a robust algebraic target. Implementing it entails real-algebraic sign determination and root-selection algorithms; merely passing a custom Hahn head into a general-purpose \code{Root} or \code{Solve} call does not install those algorithms.

\subsection{Beyond finite support: useful exact series descriptions}
Rational expressions, algebraic equations, recurrences, generating functions, and differential equations can all specify infinite series with finite data. A constructor should retain both the defining relation and the means of selecting the intended solution. Examples include
\[
 \sqrt{1+t}=1+\tfrac12t-\tfrac18t^2+\cdots,
 \qquad \exp(t)=\sum_{n\ge0}\frac{t^n}{n!},
 \qquad \sum_{n\ge0}n!t^n.
\]
All three have well-ordered support at a positive infinitesimal $t$. The last series has zero ordinary radius of convergence as a function of a complex variable; that does not prevent its strong Hahn value. It does prevent identifying it, without further choices, with an ordinary analytic germ at zero.

A useful support language should include finite supports, arithmetic progressions, positive finitely generated grids, nested lexicographic blocks, and selected well-ordered families with supplied dependency routines. More exotic computable ordinal-indexed supports can be added. They should be optional extensions, not silently treated as lists indexed by the natural numbers.

\section{When an exact description becomes an algorithm}\label{sec:effectivity}
\subsection{The two support conditions}
For a family $(a_j)$ of Hahn series, strong summability means that $\bigcup_j\supp a_j$ is well ordered and that each exponent receives contributions from only finitely many $j$. The second condition cannot be discarded. The expression $\sum_{n\ge0}1$ has support union $\{0\}$ but infinitely many contributions at exponent zero. Conversely, the support $\{1/n:n\ge1\}$ consists of positive exponents but is not well ordered. Neither expression is a strong Hahn sum.

Even an ordinary convergent scalar sum such as $\sum_{n\ge0}2^{-n}=2$ is not a \emph{strong Hahn} sum of its constant terms: exponent zero receives infinitely many contributions. Its value can instead be supplied by ordinary coefficient-field summation. The CAS must specify which summation operation is being used.

A safe constructor can recognize particular sufficient forms. The positive-support rule proves that if $E\subset\Gamma_{>0}$ is well ordered, then the additive monoid $E^*$ controls geometric and Taylor expansions. For a source support $S$, expressions indexed by one member of $S$ and a finite word in $E$ have support in $S+E^*$ and only finitely many contributions at each exponent. This is exactly the mechanism repeatedly used by AG and CT \cite{AG,CT}.

A proof of that mathematical finiteness is not automatically a method for finding the contributing words from arbitrary encodings. A CAS needs a support interface that supplies \emph{effective dependency sets}, or a presentation for which those sets can be computed. It should distinguish a proof of existence from an executable enumerator.

\subsection{A concrete finite-grid dependency algorithm}
\begin{definition}[Effective positive grid]
Fix vectors $\gamma_0,e_1,\ldots,e_m\in\Q^d$, each $e_j>0$ in lexicographic order. A grid presentation has support contained in
\[
 \gamma_0+\N e_1+\cdots+\N e_m,
\]
and supplies its coefficients either by a direct exponent procedure or by a coefficient procedure on the multi-indices, with collisions added. A finite union of such presentations is permitted.
\end{definition}

\begin{proposition}[Finite computation at a prescribed exponent]\label{prop:grid}
For an effective positive grid and a given $\lambda\in\Q^d$, all multi-indices $n\in\N^m$ with $\gamma_0+\sum_j n_je_j=\lambda$ can be computed by a finite algorithm. Consequently, Taylor or geometric constructions on finitely many such grids yield effective coefficient algorithms when their ordinary coefficient operations are effective.
\end{proposition}
\begin{proof}
There is a rational linear functional $\ell:\Q^d\to\Q$ positive on each of the finitely many $e_j$. One explicit construction starts with weight $1$ on the last coordinate and works backwards. When choosing the weight on coordinate $k$, make it sufficiently large that every $e_j$ whose first nonzero coordinate is $k$ has positive weighted sum, despite its already weighted later coordinates. Such a first nonzero coordinate is positive by lexicographic positivity. There are only finitely many inequalities.

Let $b=\ell(\lambda-\gamma_0)$. If $b<0$, there is no solution. Otherwise every solution satisfies
\[
 0\le n_j\le \left\lfloor\frac{b}{\ell(e_j)}\right\rfloor.
\]
Enumerate this finite integer box and test the original vector equation, not just its image under $\ell$. This finds exactly the contributing multi-indices. For a fixed multi-index there are finitely many ordered words and finite combinatorial multiplicities. Products and substitutions therefore give finite coefficient computations, provided the coefficient operations themselves terminate.
\end{proof}
The functional is an auxiliary finite-dependency device. It is \emph{not} an order embedding of the full higher-rank value group into $\R$, and it must not replace the lexicographic valuation elsewhere. The box algorithm is intentionally elementary; linear Diophantine solving, dynamic programming, sparse convolution, and memoized coefficient dependencies can improve its performance.

\subsection{Why arbitrary coefficient programs cannot have a zero-test}
\begin{theorem}[Representation-sensitive zero-test obstruction]\label{thm:halting}
There is no algorithm deciding whether every Hahn series presented by an arbitrary total computable rational coefficient procedure is zero. This remains true when each represented series is promised to have at most one nonzero coefficient.
\end{theorem}
\begin{proof}
Fix a positive monomial $t$ and a program $M$ on a fixed input. For $n\ge1$, define $b_M(n)=1$ if $M$ halts for the first time at step $n$, and define it to be zero otherwise. This coefficient is computable by simulating exactly $n$ steps. The expression
\[
 a_M=\sum_{n\ge1}b_M(n)t^n
\]
is a valid Hahn series with support either empty or a singleton. It is zero precisely when $M$ never halts. A total zero-test for these coefficient presentations would decide the halting problem, which is impossible.
\end{proof}
Every value in this example is actually a polynomial monomial or zero. There is no contradiction with \cref{thm:rationalcore}: an explicit finite polynomial list is different information from a program that might reveal one nonzero coefficient arbitrarily late. The obstruction concerns the representation, not just a set of mathematical values.

A leading-term algorithm that always either reports zero or returns the least nonzero exponent would also decide zero, so it cannot exist for these arbitrary programs. Numerical sampling does not repair the problem. Returning \code{False} after a fixed unsuccessful search is unsound; the correct result is an explicit unknown status or a resource-limited inconclusive result.

\subsection{Positive zero-testing results must also be retained}
This obstruction does not rule out rich structured languages. Rational and algebraic representations already avoid it. For an ordinary-point solution of a scalar linear differential equation with analytic coefficients and nonvanishing highest-order coefficient, finitely many initial derivatives determine the solution. With effective coefficient arithmetic and the equation recorded, vanishing of the requisite initial data is a finite zero certificate. Singular equations and arbitrary differential-algebraic solution descriptors require more careful selection data.

A particularly relevant current result is Chen--Fang--van der Hoeven's \emph{A Zero-Test for D-Algebraic Transseries}, published in the ISSAC 2026 proceedings \cite{CFH}. Their framework uses effective differential ambient fields and structured grid-based transseries. Theorem 18 extends the ambient field, and its zero-test, by a distinguished solution of an appropriate quasi-linear differential equation. Ritt reduction and separation bounds are central. This is a concrete route beyond elementary expressions, not a zero-test for arbitrary series-producing programs. A CAS implementing it must preserve its effective-field, transbasis, and solution-selection assumptions rather than replacing them with an unrestricted \code{DifferentialRoot} label.

\section{Precision without loss of valuation rank}\label{sec:precision}
\subsection{A finite truncation need not reach a requested scale}
Take $\Gamma=\Q\omega+\Q$ and put $u=t^\omega$. Then $u$ is smaller than $t^n$ for every ordinary positive integer $n$. The exact remainder identity is
\begin{equation}\label{eq:geometricremainder}
 \frac1{1-t}-\sum_{j=0}^{N-1}t^j=\frac{t^N}{1-t},
 \qquad v\left(\frac{t^N}{1-t}\right)=N<\omega.
\end{equation}
No finite prefix of this expansion gives error valuation at least $\omega$. Yet the strong sum exists, and the rational expression on the left is a finite exact description. The supplied manuscripts emphasize this difference between strong sums and sequential valuation convergence \cite{AG,CT}.

The same issue affects $\sum_{n\ge0}t^n+u$: an ordinary stream that insists on outputting all terms in increasing valuation order will never reach the $u$ term. A suitable representation can instead store a whole $t$-block as a rational or lazy coefficient in an outer $u$-series. Another suitable representation provides direct exponent access. Neither mechanism is equivalent to taking more digits in a single real-valued precision parameter.

Even in rank one, an arbitrary Hahn support need not have finite intersection with a bounded interval. The set $\{1-1/n:n\ge1\}$ is increasing and well ordered, but has infinitely many exponents below $1$. Positive finite grids in an Archimedean finitely presented setting are a useful restricted alternative. A request for ``all terms below a threshold'' needs a finiteness guarantee or must return a compressed infinite block.

\subsection{Certified error objects}
Use the convention
\[
 x=\widetilde x+O_v(\alpha)\quad\Longleftrightarrow\quad
 v(x-\widetilde x)\ge\alpha.
\]
The precision parameter is an element of the full ordered group. It is not an ordinary numerical accuracy unless an additional interpretation has been chosen. For $v(x-\widetilde x)\ge\alpha$ and $v(y-\widetilde y)\ge\beta$, one has
\begin{align}
 x+y&=\widetilde x+\widetilde y+O_v(\min(\alpha,\beta)),\label{eq:addprecision}\\
 xy&=\widetilde x\widetilde y+
 O_v\!\left(\min\{\alpha+v(\widetilde y),\beta+v(\widetilde x),
 \alpha+\beta\}\right).\label{eq:mulprecision}
\end{align}
These follow by expanding the errors and applying the valuation inequality. If $\alpha>v(\widetilde x)$, then $x$ and $\widetilde x$ have the same valuation, and
\begin{equation}\label{eq:invprecision}
 x^{-1}=\widetilde x^{-1}+O_v\bigl(\alpha-2v(\widetilde x)\bigr).
\end{equation}
Indeed, $x^{-1}-\widetilde x^{-1}=-(x-\widetilde x)/(x\widetilde x)$.

For complex coefficients, an additional ordinary coefficient-accuracy layer may be used, but should remain separate. An error in the leading ordinary coefficient cannot be dismissed as a higher Hahn-order error. Certified interval or ball arithmetic for coefficients and exact valuation truncation solve different problems.

\subsection{Jets, nested precision, and exact tails}
A finite Taylor jet in $K[z]/(z)^N$ has exact quotient-ring semantics. It determines products modulo $(z)^N$, not a unique analytic function. An asymptotic object has a different meaning: it records a function or value modulo an explicitly stated error class. An exact lazy object specifies all coefficients, even though only finitely many are cached. An exact rational tail can specify infinitely many coefficients in one node. These distinctions should survive serialization and printing.

Wolfram's \code{SeriesData} already demonstrates the usefulness of a separate series object carrying an expansion point, rational exponent step, and omitted-order information \cite{WLSeries}. It can serve as a backend for particular one-parameter jets. It is not by itself an arbitrary support language for $\Q^d_{\mathrm{lex}}$, nor does \code{Normal} preserve the omitted tail. Nested \code{SeriesData} objects can be useful when a fixed iterated expansion convention is supplied; their existence does not automatically define a general multi-rank Hahn truncation policy.

\section{Which functions admit exact symbolic evaluation?}\label{sec:functions}
\subsection{Algebraic functions and elementary field operations}
Polynomials and rational functions evaluate exactly at every represented tuple where the denominators are proved nonzero. Integer powers stay in the field. Rational powers generally require algebraic extension and root selection. Conjugation, real and imaginary parts, and squared modulus are finite field operations. Modulus is
\[
 |a+ib|=\sqrt{a^2+b^2},
\]
using the nonnegative root in a real closed extension. The minimal rational-monomial package therefore implements squared modulus but not unrestricted modulus.

Real comparison gives \code{Min}, \code{Max}, and piecewise evaluation when all required signs are decidable. A complex field has no compatible order. For an infinite surreal, an ordinary $\Z$-valued floor cannot satisfy the usual defining inequalities because no ordinary integer brackets that value. A package offering rounding to surreal or omnific integers must specify that different target and its rule, rather than silently reusing ordinary integer rounding.

Polynomial equations and inequalities over an effective real algebraic extension can be attacked by finite sign algorithms, Sturm methods, and real-algebraic elimination. An arbitrary surreal constant represented only by a symbol does not magically supply the comparison oracle these algorithms require. Quantifier elimination for a mathematical field and executable elimination over a chosen input encoding are separate assertions.

\subsection{Canonical infinitesimal Taylor lifts}
Let $f$ be an ordinary holomorphic function near $c\in\C$, with a specified branch where necessary. For an infinitesimal $\epsilon$, define
\begin{equation}\label{eq:Taylorlift}
 f^\sharp(c+\epsilon)=
 \sum_{n\ge0}\frac{f^{(n)}(c)}{n!}\epsilon^n.
\end{equation}
If the input and coefficient procedures have effective support data, \cref{prop:grid} turns this into a coefficient algorithm. The exact value is a lazy expression, not merely any displayed finite Taylor polynomial. The ordinary coefficients may enlarge the coefficient parent.

For example,
\begin{align}
 \exp^\sharp(c+\epsilon)&=e^c\sum_{n\ge0}\epsilon^n/n!,\label{eq:localexp}\\
 \log^\sharp(c+\epsilon)&=\log(c)+
 \sum_{n\ge1}\frac{(-1)^{n+1}}{n}(\epsilon/c)^n,
 \qquad c\ne0,\label{eq:locallog}\\
 \sin^\sharp(c+\epsilon)&=\sin(c)\cos^\sharp(\epsilon)
 +\cos(c)\sin^\sharp(\epsilon).\label{eq:localsin}
\end{align}
Here the logarithm's ordinary germ must be chosen; there is no holomorphic logarithm germ at zero. At $c=0$, the exponential coefficients lie in $\Q$, whereas at an algebraic nonzero $c$ the coefficient $e^c$ may require a transcendental constant descriptor. Domain closure must be checked at both the coefficient and exponent levels.

\begin{proposition}[Correctness of infinitesimal substitution]\label{prop:taylor}
Suppose $f$ has a specified ordinary analytic germ at $c$ and $\epsilon$ has well-ordered positive support $E$. Then \eqref{eq:Taylorlift} is a strong Hahn sum with support contained in $E^*$. Sums, products, and compositions of such analytic germs are respected whenever the ordinary centers match and the relevant germs are defined.
\end{proposition}
\begin{proof}
Every coefficient contribution is indexed by a finite word in $E$; the positive-support lemma gives strong summability. At a fixed exponent, the formal Taylor identities for sums, products, and composition involve only finitely many contributing terms. They therefore remain valid after substitution. Center values are handled as ordinary coefficients, before the infinitesimal expansion is performed.
\end{proof}
This is the local construction used by SA. It does not extrapolate an ordinary function to every infinite argument. For $\epsilon\ne0$ infinitesimal, the proposed substitution
\[
 \sum_{n\ge0}\frac{\epsilon^{-n}}{n!}
\]
has strictly descending valuations. It is not a strong Hahn definition of $e^{1/\epsilon}$. A separate global real exponential or surcomplex extension may assign a value, but that is a different construction \cite{SA}.

\subsection{Special functions and differential descriptions}
Many useful special-function germs can be specified by a finite differential equation or recurrence, initial values, and a chosen ordinary domain. A linear differential equation at an ordinary point provides recursive Taylor coefficients; algebraic equations provide implicit-function recurrences at simple branches. Parameter singularities, resonance, irregular singular points, and branch changes must be part of the input checks.

Thus a useful \code{AnalyticGerm} descriptor contains the center, ordinary function or differential equation, coefficient parent, branch identifier, coefficient procedure, and evidence that the chosen ordinary germ exists. Calling \code{SeriesCoefficient} on a supported classical function can supply coefficients, but a returned unevaluated expression is not a termination guarantee for all requested indices. User-supplied recurrences require their own totality and domain assumptions.

A general smooth function does not have a unique infinitesimal extension determined by its Taylor coefficients. For example, an ordinary smooth function flat at zero has the same Taylor jet as the zero function there. Assigning only that jet to all infinitesimals discards information about its nonzero values away from zero. CT explicitly separates analytic target coefficients from merely smooth parameter coefficients; a CAS should retain the same distinction \cite{CT}.

\subsection{Exact formal values versus analytic and resurgent values}
The series $F(t)=\sum_{n\ge0}n!t^n$ is an exact strong Hahn value at an infinitesimal $t$, and its coefficients satisfy a simple recurrence. As an ordinary complex power series it is divergent. In the formal parameter sense it satisfies
\[
 t^2 F'(t)+(t-1)F(t)=-1.
\]
This identity can be checked coefficientwise: at positive degree $n$, the contribution is $n(n-1)!-n!=0$, and the constant contribution is $-1$. An ordinary differential equation with such a formal asymptotic solution may have different analytic solutions sharing the same expansion, with differences invisible to every power of $t$.

Consequently, an exact formal Hahn descriptor and an exact analytic-function descriptor should remain distinct. A Borel, sectorial, or resurgent summation policy may relate them, but must record directions, singularities, and normalization data. Costin--Ehrlich provide substantial surreal extension and integration constructions for suitable analytic and resurgent classes \cite{CE}; this is relevant infrastructure, not a license to identify every formal asymptotic series with a unique global surcomplex function.

\section{Exponentials, transseries, and noncanonical surcomplex phases}\label{sec:transseries}
\subsection{Real Gonshor exponential and changing the monomial domain}
The established surreal exponential $\exp_{\No}$ extends the real exponential and is an ordered-group isomorphism from $(\No,+)$ to $(\No_{>0},\cdot)$. Its inverse is $\log_{\No}$ \cite{Gonshor}. The identities
\[
 \exp_{\No}(x+\epsilon)=\exp_{\No}(x)\sum_{n\ge0}\epsilon^n/n!
\]
for real infinitesimal $\epsilon$ are compatible with local Taylor lifting. They do not say that the normal form of $\exp_{\No}(x)$ can be computed by the ordinary exponential Taylor series at an arbitrary infinite $x$.

A finite expression such as $\exp_{\No}(\omega)$ is an exact name once that operation is part of the language. For an effective implementation, the exponential should decompose the argument according to its large, constant, and small parts; apply the small-part series where valid; and introduce or recognize a new growth monomial for the large part. Normalization, comparison, and logarithmic inversion then require a transseries or surreal-exponential backend, not just a new display head.

A fixed Hahn workspace containing the input may not be closed under these transcendental operations. Closure under algebraic roots in a divisible complex Hahn field is a different issue. Bournez--Guilmant study set-sized surreal fields stable under exponential, logarithmic, derivative, and antiderivative operations \cite{BG}. Their structural closure conditions help choose mathematical domains; a bound on birthdays is not by itself a finite executable representation or a complexity estimate.

\subsection{Transseries as a practical middle layer}
A transseries backend represents hierarchies involving powers, logarithms, and exponentials, with dominant monomials and controlled support. It is much closer to symbolic asymptotic computation than a raw tree of Conway options. Van der Hoeven's monograph develops the pertinent algebraic and computational framework \cite{vdH}. The 2026 zero-test discussed in \cref{sec:effectivity} gives an important additional algorithmic direction for structured differential-algebraic extensions \cite{CFH}.

Three pieces should be recorded separately: the abstract transseries expression, the chosen realization in a surreal field, and the permitted function-composition or differential structure. Berarducci--Mantova construct surreal derivations and study transseries as surreal function germs \cite{BMder,BMfun}. These results support substantial exact function languages, but do not identify every arbitrary Hahn series with a globally composable function on all of $\No$.

For a first implementation, a finite-depth exponential--logarithmic tower with explicit dominant-term certificates is a reasonable design target. Its scope should be stated by grammar and algorithms, not by saying it ``supports all transseries''. The broader differential-algebraic algorithms are a separate implementation project. Terms outside the supported class can remain exact formal names with reduced capabilities, provided their mathematical meaning has actually been fixed.

\subsection{Why the surcomplex exponential needs a policy}
At finite imaginary arguments, ordinary sine and cosine have canonical infinitesimal Taylor lifts. At infinite imaginary arguments, the supplied SA manuscript shows that natural-looking local requirements do not determine one exponential. Let $\operatorname{fin}(y)$ retain the nonnegative-exponent part of a real surreal normal form, and let $\ell(y)$ be the coefficient of $t^{-1}=\omega$. For $\alpha\in\R$, SA constructs
\begin{equation}\label{eq:phaseexp}
 E_\alpha(x+iy)=\exp_{\No}(x)
 \exp_{\mathrm{fin}}\bigl(i(\operatorname{fin}(y)+\alpha\ell(y))\bigr).
\end{equation}
These are fine-locally analytic exponential group homomorphisms extending the ordinary complex exponential, with derivative equal to themselves. Nevertheless,
\begin{equation}\label{eq:phasecontrast}
 E_0(i\omega)=1,\qquad E_\pi(i\omega)=-1.
\end{equation}
We use this as a \emph{manuscript result and an interface warning}, not as a published uniqueness classification \cite[Section 11]{SA}.

A proposed \code{SurcomplexExp[z]} must therefore either restrict its input domain, specify a particular extension policy, or remain unassigned at such arguments. A policy identifier belongs in the parent or function descriptor and in cache keys. Discarding it would allow the same stored argument to evaluate to incompatible values. The same warning applies to unrestricted trigonometric functions and to expressions such as $E_\alpha(1/z)$ near a microscopic essential singularity.

The policy is only useful computationally when operations such as extracting $\operatorname{fin}(y)$ and $\ell(y)$ are effective for the input representation. They are mathematical operations on normal forms, but need not be computable by scanning an arbitrary coefficient stream. The representation problem therefore reappears even after the semantic choice has been settled.

\subsection{Logarithms, powers, and branch-sensitive simplification}
The ordinary logarithm identities require branch checks, and the surcomplex setting adds extension choices. The local identity
\[
 \log(1+\epsilon)+\log(1+\eta)
 =\log\bigl((1+\epsilon)(1+\eta)\bigr)
\]
is valid for the infinitesimal logarithm and infinitesimal $\epsilon,\eta$. A global identity $\log(zw)=\log z+\log w$ is not justified by this local calculation. Likewise, \code{PowerExpand}-style rewrites cannot be applied without recording the relevant branch and domain assumptions.

SA constructs a global fine-analytic logarithm after choosing the ordinary argument of the standard unit direction. It is a right inverse to its specified exponentials and has derivative $1/z$, but it is not a Hahn-coherent primitive on the punctured ordinary base \cite{SA}. Thus a CAS must not combine that logarithm with a coherent contour fundamental theorem to deduce $\oint dz/z=0$. The function category is part of the theorem's type.

SA further proves, under its explicit all-scale coherence definition, that functions coherent entire at \emph{every surreal radius} are polynomials, and the analogous meromorphic functions are rational. This manuscript rigidity statement explains why requiring one universal all-scale analytic class is not a practical way to include global transcendental functions. A smaller admissible atlas or a different function framework is needed; the article does not claim that any one such global atlas has already been implemented.

\section{Differentiation, germs, and finite analytic computation}\label{sec:germs}
\subsection{Three derivatives that must not share an implicit convention}
The derivatives in the supplied analytic manuscripts are derivatives in ordinary or surcomplex variables $z_j$, with the Hahn scalar field $K$ held constant. Thus
\[
 D_z\bigl(t^\gamma f(z)\bigr)=t^\gamma f'(z).
\]
Coefficientwise differentiation preserves the containing Hahn support. It is the appropriate operation for local holomorphic identities, Jacobians, residues, and the contour calculus of CT \cite{SA,AG,CT}.

A formal parameter derivative is different. In a rational or Puiseux field with $t$ as an independent parameter and constant rational exponent $q$, it satisfies $d(t^q)/dt=qt^{q-1}$. It changes valuations and may decrease precision. In a multivariate rational-monomial field, each formal partial derivative treats the other independent variables as constants.

A surreal derivation is additional global structure on $\No$. The Berarducci--Mantova construction supplies a distinguished differential-field framework compatible with exponential structure \cite{BMder}. It should be invoked through a separate operator and parent. For instance, after choosing a derivation with $\partial\omega=1$, one obtains $\partial(\omega^{-1})=-\omega^{-2}$; by contrast $D_z(\omega^{-1})=0$. Applying a constant-exponent rule when the exponent itself is a nonconstant surreal for that derivation can be wrong. A default \code{D} overload cannot infer all these intentions from a displayed monomial.

\subsection{The radius-free Hahn-germ algebra}
AG uses
\begin{equation}\label{eq:germalgebra}
 R_n=\C\{z_1,\ldots,z_n\},\qquad
 \Acal_n=R_n((t^\Gamma)).
\end{equation}
Each coefficient is an ordinary convergent germ at zero, but different coefficients need not have a common ordinary radius. A finite computational descriptor should therefore separate its Hahn-support program from its ordinary coefficient-germ program. A common-domain family in $\OO(U)((t^\Gamma))$ is a stronger and different certificate \cite{AG}.

The distinction has executable consequences. Let
\begin{equation}\label{eq:radiusfree}
 H(z)=\sum_{n\ge1}\frac{t^n}{1-nz}.
\end{equation}
Every coefficient is a valid ordinary germ at zero. There is no positive ordinary disk on which all coefficients are holomorphic, since their poles are $1/n$. Nevertheless $H(a)$ is well defined for every infinitesimal $a$, using strong Taylor substitution. Calling an ordinary contour routine on a fixed disk containing all coefficient domains would be an error; infinitesimal rescaling provides the appropriate alternative in CT.

AG also distinguishes \eqref{eq:germalgebra} from $K[[z]]$. The formal expression
\[
 (z-t^\eta)^{-1}=-\sum_{j\ge0}t^{-(j+1)\eta}z^j
\]
is a power series over $K$ but not an element of $\Acal_1$: its combined Hahn exponents are not well ordered. In $\Acal_1$, the function $z-t^\eta$ actually vanishes at an infinitesimal point, so it must not be marked as an invertible germ on the whole monad. A CAS needs separate types for a formal completion at a point and a function algebra on the entire infinitesimal neighborhood \cite{AG}.

SA also permits sufficiently small full-surreal rescalings of much more general formal local series \cite[Section 4]{SA}. The required scale may lie outside a previously chosen workspace. An existence theorem for such a scale is not a uniform routine for finding it from arbitrary coefficient programs; an effective implementation needs bounds or a construction supported by its input language.

\subsection{Evaluation, translation, and finite jets}
For $F=\sum_\gamma f_\gamma t^\gamma$ and an infinitesimal tuple $a$, AG defines
\begin{align}
 F(a)&=\sum_{\gamma,\alpha}
 t^\gamma\frac{\partial^\alpha f_\gamma(0)}{\alpha!}a^\alpha,\label{eq:germeval}\\
 (\tau_aF)(z)&=\sum_{\gamma,\alpha}
 t^\gamma\frac{\partial^\alpha f_\gamma(z)}{\alpha!}a^\alpha.\label{eq:germtranslation}
\end{align}
The support certificate is $S+E_a^*$, where $S=\supp F$ and $E_a$ is the union of the coordinate supports of $a$. A coefficient request can be compiled into a finite dependency graph when those supports have effective presentations and the ordinary germ derivative requests terminate. This gives an actual algorithmic interpretation of AG's formula, not an assumption that every convergent germ is already a CAS object.

Finite jets are particularly accessible. At a represented infinitesimal $a$, the manuscript identifies
\begin{equation}\label{eq:finitejets}
 \Acal_n/(z-a)^N\simeq
 K[X_1,\ldots,X_n]/(X_1,\ldots,X_n)^N.
\end{equation}
After the required finite derivatives have been computed, the right-hand side uses ordinary sparse polynomial arithmetic with a degree cutoff. A separate Hahn precision cutoff may be applied to its coefficients. The Taylor degree and the Hahn valuation are different axes of approximation and should not be compressed into one integer option.

For a formal implicit system with invertible Jacobian, coefficients can be computed recursively using linear solves over the scalar parent. A returned solution should include the root center, leading branch data, Jacobian certificate, and support or precision bounds. A singular Jacobian is not a proof that no solution exists; it means that this simple-root procedure is inapplicable and a preparation, ramification, or finite-algebra method is needed.

\subsection{Preparation and division as operator algorithms}
AG's one-variable preparation and its several-variable division formulas are useful implementation specifications. For $F=f+E$, with leading ordinary germ $f$ regular in a selected variable $w$ and $v_H(E)>0$, let $D$ and $\Pi$ be ordinary Weierstrass quotient and remainder operators. Then
\begin{equation}\label{eq:weierstrassalgorithm}
 Q=(\id+D M_E)^{-1}DG
 =\sum_{r\ge0}(-D M_E)^rDG,\qquad
 R=\Pi(G-EQ).
\end{equation}
The support is contained in $\supp G+(\supp E)^*$. These are exact formulas from AG, not an implemented routine in the accompanying small package \cite{AG}.

To execute them, the system needs a concrete ordinary division backend. For polynomial or algebraic coefficient germs, polynomial reduction and selected local algebra procedures are natural choices. For a suitable differential or analytic descriptor, additional constructive division methods may be possible. An arbitrary complex-linear splitting chosen by the axiom of choice is not executable merely because it preserves the space of analytic germs. Even a finite output coefficient can call an ordinary operation that is unavailable for the given germ language.

The ordinary leading order and regularity direction must also be found or supplied. If detecting whether a coefficient germ is zero is only semidecidable, searching for the first nonzero homogeneous part need not terminate on zero input. A reliable API accepts certified leading data, tries supported methods, and otherwise returns a structured unresolved result rather than asserting a Weierstrass polynomial.

\section{Finite zero schemes, normal forms, and residues}\label{sec:finitealgebra}
\subsection{A computationally rich target beyond scalar arithmetic}
For a leading isolated complete intersection
\[
 f=(f_1,\ldots,f_n),\qquad
 \mu=\dim_\C R_n/(f_1,\ldots,f_n)<\infty,
\]
and a positive-support perturbation $F=f+E$, AG proves preservation of the finite algebra length and supplies multiplication matrices, residues, and trace identities \cite{AG}. This is an especially attractive CAS target because the result can be represented by finitely many matrices even when their entries are exact lazy Hahn series.

Choose a monomial basis $b_1,\ldots,b_\mu$ of the leading quotient. An ordinary splitting is written
\[
 G=i p(G)+\sum_j f_jh_j(G),\qquad pi=\id,\qquad hi=0.
\]
For an effective implementation, $p$ and $h_j$ must be algorithms for the chosen coefficient-germ presentation. AG's arbitrary existence of a linear splitting is sufficient for its theorem but not for a program.

With $T=\sum_j E_jh_j$, the manuscript's normal-form coordinates are
\begin{equation}\label{eq:NFalgorithm}
 \NF_F(G)=p(\id+T)^{-1}G
 =\sum_{r\ge0}(-1)^r pT^rG.
\end{equation}
The matrices of multiplication are
\begin{equation}\label{eq:multiplication}
 (M_H)_{rs}=\NF_F(Hb_s)_r.
\end{equation}
A finite-grid support backend and a constructive ordinary splitting give effective coefficient access to every entry. Matrix multiplication, determinant, trace, and characteristic polynomial then reduce to the corresponding scalar operations. None requires explicitly isolating every perturbed root.

\subsection{Residue computation and finite certificates}
AG defines the perturbed residue by
\begin{equation}\label{eq:residuealgorithm}
 \Res_F(H)=\sum_{k\in\N^n}(-1)^{|k|}
 \res_{f^{k+1}}(H E_1^{k_1}\cdots E_n^{k_n}).
\end{equation}
Its support lies in $\supp H+E^*$. At a requested exponent, an effective dependency routine reduces the calculation to finitely many ordinary residues. When the leading denominators are coordinate powers, those ordinary residues are simply Taylor coefficient extraction. More general finite-dimensional residues have established algebraic computation methods \cite{CDS}.

The manuscript trace identity is
\begin{equation}\label{eq:traceidentity}
 \Res_F(HJ_F)=\Tr(M_H),\qquad
 J_F=\det\left(\frac{\partial F_i}{\partial z_j}\right).
\end{equation}
This gives both a useful computation and a consistency check. For $H=1$, the answer is the ordinary integer $\mu$. A finite-algebra object should record basis, multiplication matrices, defining relations, a residue vector or pairing matrix, and enough reduction data to verify the relations. Such a representation retains multiplicities and nilpotents that a list of distinct points would lose.

The hypotheses matter computationally. AG's length-preservation theorem concerns $n$ equations in $n$ variables with an isolated leading complete intersection. An arbitrary overdetermined system is not covered. Its example $(z^2,0)\mapsto(z^2,t^\eta z)$ changes length from two to one. An implementation must test or demand the isolated-complete-intersection certificate, not apply a conservation rule to every infinitesimal perturbation \cite{AG}.

\subsection{Root separation and certified continuation}
For a simple zero $a$ of a nonnegative-support system, AG sets
\[
 \kappa_a=\max\left(0,-\min_{i,j}v\bigl((DF(a)^{-1})_{ij}\bigr)\right).
\]
For a perturbation $D$ with $\delta=\min_i v_H(D_i)>2\kappa_a$, it proves a unique continued zero in the ball $v_{\min}(b-a)>\kappa_a$ and the displacement bound $v_{\min}(b-a)\ge\delta-\kappa_a$ \cite{AG}. These are exact inequalities in the full ordered value group.

A root tracker can use this result as a certificate-producing operation: return the old root descriptor, the perturbation bound, the inverse-Jacobian bound, and the new implicit root with a verified uniqueness ball. It need not enumerate an infinite Newton sequence. The equality case cannot be accepted: $z^2-t^{2\eta}$ perturbed by $t^{2\eta}$ collapses two simple roots into a double root precisely at $\delta=2\kappa_a$.

These guarantees remain conditional on obtaining the exact bounds. A numerical estimate at one real specialization of the monomials is not a proof of the higher-rank inequality. Interval bounds for ordinary coefficients can assist a computation, but cannot replace exponent-group comparisons.

\section{Integration and contour operations}\label{sec:integration}
\subsection{Primitives in a specified differential algebra}
Rational integration in a formal variable $z$ over a characteristic-zero coefficient field is a well-defined algebraic problem. Partial fractions show that polynomial terms and higher-order poles have rational primitives, while a simple pole contributes a logarithmic obstruction. Hence a rational differential has a rational primitive exactly when all its simple-pole residues vanish, after passing to an algebraic splitting field as necessary.

For common-domain Hahn families $G=\sum_\lambda g_\lambda t^\lambda$ on an ordinary domain $U$, SA defines differentiation and integration coefficientwise. A coherent primitive exists precisely when its ordinary closed-cycle Hahn periods vanish. The same manuscript exhibits fine-locally constant functions with zero derivative but different values on different fine-open pieces. There can therefore be no unrestricted endpoint integral satisfying the fundamental theorem for every fine-locally analytic function \cite{SA,CT}.

A proposed integration API should distinguish rational primitive, analytic-germ primitive, coherent primitive on a stated domain, and antiderivative for a specified surreal derivation. A successful result in one category is not a result in every other category. Additive constants can be normalized locally, but a global primitive also needs compatibility across the relevant domain.

\subsection{Coherent and algebraic closed-contour functionals}
For an ordinary piecewise smooth path $\gamma$ and common-domain Hahn data, CT uses
\begin{equation}\label{eq:Hintegral}
 \HInt_\gamma G(z)\,dz
 =\sum_\lambda t^\lambda\int_\gamma g_\lambda(z)\,dz.
\end{equation}
This is a coefficientwise definition on an ordinary parameterized chain. It is not a limit of arbitrary fine-topological Riemann sums. An exact CAS operation can return a lazy series of ordinary integrals when those integrals are represented, or use Cauchy and residue rules to avoid integration altogether \cite{CT}.

For a rational differential and a finite semialgebraic cycle avoiding its poles, CT also defines an algebraic period by
\begin{equation}\label{eq:algebraicperiod}
 I_C^{\mathrm{alg}}(R\,dz)
 =2\pi i\sum_{a\text{ a finite pole}} w_F(C,a)\res_a(R\,dz).
\end{equation}
The winding number can be obtained from certified ray crossings or appropriate real-algebraic index computations. This is an especially practical arbitrary-scale operation when the scalar ordered field is effective. The normalized period $I_C^{\mathrm{alg}}/(2\pi i)$ may remain in an algebraic coefficient domain even when multiplying by $2\pi i$ requires adjoining a named transcendental constant. CT states agreement with Hahn integration on their common admissible contours \cite{CT}.

The data type of a contour must include orientation and pole-avoidance conditions. A pole on the contour is not a zero winding number; it is outside the definition unless a separate regularization has been selected. The system should not infer a Cauchy theorem merely from the fact that a path is closed: null-homology or an appropriate residue calculation is necessary.

\subsection{Microscopic circles and separating tori}
CT rescales a radius-free germ by $z=a+t^s w$, with infinitesimal $a$ and positive $s$. At each Hahn exponent, the resulting coefficient in $w$ is an ordinary polynomial, because only finitely many Taylor degrees contribute. This supplies a common ordinary domain \emph{after} rescaling, even when none existed before. The support certificate, rather than a guessed radius, authorizes microscopic Cauchy and residue operations \cite{CT}.

For the finite algebra in \cref{sec:finitealgebra}, form coordinate characteristic polynomials
\[
 Q_i(T)=\det(TI-M_{z_i})=T^\mu+\sum_{j<\mu}q_{ij}T^j.
\]
Using support-controlled normal forms, CT obtains a matrix $\mathsf C$ such that $Q_i(z_i)=\sum_j C_{ij}F_j$. Choose positive $s_i$ satisfying
\begin{equation}\label{eq:radiusineq}
 (\mu-j)s_i<v(q_{ij})\quad\text{for each nonzero }q_{ij}.
\end{equation}
With $r_i=t^{s_i}$, its separating-torus formula is
\begin{equation}\label{eq:torusformula}
 \Res_F(H)=\frac{1}{(2\pi i)^n}
 \HInt_{\mathbb T_r}
 \frac{H\det\mathsf C}{\prod_i Q_i(z_i)}
 \,dz_1\wedge\cdots\wedge dz_n.
\end{equation}
Here the torus is parameterized by $z_i=r_i e^{i\theta_i}$ using ordinary angular parameters. Formula \eqref{eq:torusformula} is imported conditionally from CT's stated finite-algebra and transformation hypotheses \cite{CT}.

This yields a concrete compilation pipeline: compute multiplication matrices, characteristic polynomials, an ideal-membership transformation matrix, and valuation inequalities; then reduce the integral to coefficient extraction. The determinant must be retained. The original denominators $F_i$ need not avoid the chosen product torus, and a real $n$-torus is not the whole real $(2n-1)$-dimensional boundary of a polydisk. These are mathematical conditions that a symbolic simplifier can otherwise accidentally erase.

\subsection{Stokes as a symbolic verification rule}
CT's Hahn differential forms have ordinary smooth coefficient forms with a common well-ordered support. Exterior derivative, wedge product, admissible pullback, and integration along compact ordinary parameter chains satisfy the coefficientwise Stokes identity. Positive-support deformations enlarge the class of chains under stated analyticity conditions \cite{CT}.

A CAS can use this theorem as a certificate-based rewrite: replace the integral of $d\omega$ over a certified chain by the integral of $\omega$ over its oriented boundary, or prove invariance under a certified pole-avoiding homotopy. This is a proposed use of an identified theorem, not an algorithm deciding whether every arbitrary input map is admissible. Nor does definable Jordan separation alone furnish a general field-valued integration operation.

\section{A Wolfram Language architecture}\label{sec:architecture}
\subsection{Parent domains, representations, and capabilities}
An implementation should have a small trusted algebraic core and separate extension modules. A parent records coefficient domain, exponent domain, monomial embedding, support presentation, and available algorithms. A value records its parent and representation: sparse finite form, normalized fraction, selected algebraic root, lazy grid, or function-derived expression. A function records its variables, coefficient parent, ordinary or nonstandard domain, regularity category, branch, and extension policy.

A schematic interface might be:
\par\noindent\begin{minipage}{\textwidth}
\begin{lstlisting}
HahnParent[coefficientDomain, exponentDomain,
  "MonomialInterpretation" -> embedding,
  "SupportClass" -> supportGrammar]
HahnValue[parent, representation, certificates]
AnalyticGerm[center, definition, initialData, domain, branch]
CoherentFamily[parent, ordinaryDomain, coefficientProgram]
SurcomplexFunction[definition, domain, extensionPolicy]
\end{lstlisting}
\end{minipage}\par
These are design sketches, not symbols implemented by Wolfram Language or by the small package. The constructor should state whether its evidence has been checked, supplied as an assumption, or inherited from a proved construction. User assertions must not be relabeled as kernel-verified theorems.

Capabilities should include arithmetic, exact equality, real comparison, leading term, coefficient access, finite-cut truncation, compressed truncation, algebraic root selection, composition, and specified differentiation. Different values can expose different subsets. A general symbolic layer can still compose expressions, but cannot claim the algorithmic guarantees of its strongest subdomain for every expression it contains.

\subsection{Reusing existing facilities without changing their meaning}
\begin{center}\small
\begin{longtable}{@{}L{0.23\textwidth}L{0.33\textwidth}L{0.35\textwidth}@{}}
\toprule Wolfram facility & Appropriate use & Extra surreal-specific work\\\midrule\endhead
\code{Together}, \code{Cancel}, polynomial routines & Exact fraction and finite polynomial arithmetic & Declare independent monomials and their ordered exponent parent.\\[0.35em]
\code{Root}, \code{AlgebraicNumber} & Ordinary exact algebraic coefficients & Supply root selection over non-Archimedean coefficient fields; do not borrow ordinary complex ordering blindly.\\[0.35em]
\code{SeriesData}, \code{SeriesCoefficient} & Supported ordinary or one-parameter Taylor/Puiseux calculations & Strong-support certificates, higher-rank exponents, and exact-tail semantics.\\[0.35em]
\code{Asymptotic}, related functions & Propose or compute asymptotic expansions & Prove the selected expansion's surreal interpretation; retain beyond-all-orders and branch data.\\[0.35em]
\code{D}, algebraic matrices & Ordinary coefficient derivatives and finite algebra calculations & Identify which variables and monomials are constant for the chosen derivation.\\[0.35em]
\code{Inactive}, custom heads & Preserve exact unevaluated descriptions & A denotation, domain checks, and algorithms are still required.\\[0.35em]
\code{UpValues}, formatting rules & Local arithmetic dispatch and readable display & Avoid broad rewrites, implicit parent mixing, or replacing semantic data by boxes.\\
\bottomrule
\end{longtable}\end{center}
The documented contracts for these facilities provide useful building blocks, not an automatically installed arbitrary surreal field \cite{WLRoot,WLAlgebraic,WLSeries,WLAsymptotic,WLInactive,WLUpValues}. This is a statement about the examined interfaces, not a claim to have exhaustively surveyed every third-party Wolfram package.

The same distinction is visible in other CAS designs. Sage's lazy-series documentation explicitly distinguishes coefficients computed on demand from equality, which is decidable only in special cases \cite{SageLazy}. Its asymptotic rings separate coefficient rings, growth groups, exact terms, and order terms \cite{SageAsymptotic}. These are useful architectural precedents. Their existing domain definitions should be reused where they fit, not relabeled wholesale as the full surreal class.

\subsection{Evaluation control and safe dispatch}
A first prototype should use explicit functions such as \code{HAdd}, \code{HMul}, and \code{HCompare}. Carefully scoped upvalues may later make familiar arithmetic convenient, but the package should never broadly unprotect or redefine \code{Plus}, \code{Times}, \code{Power}, \code{Exp}, or \code{Less}. The official upvalue mechanism allows local definitions associated with custom objects \cite{WLUpValues}; it does not remove the responsibility to check mixed parents and analytic hypotheses.

Keep raw monomial symbols internal or require them to be unassigned. Ordinary substitution can destroy their intended relationships: replacing $u=t^\omega$ by $t^{100}$ collapses a rank. Promoting two parents requires an explicit embedding preserving coefficient arithmetic, exponent relations, order, and function policy. A different variable order is a different parent, even if the printed rational expression is identical.

Results should distinguish invalid input, undefined mathematical operation, unsupported algorithm, and unknown proposition. A failed zero-test must not return \code{False}; a negative-power operation at zero is genuinely undefined; a square root outside a rational parent asks for an extension rather than proving impossibility. On general symbolic inputs, a conditional result should preserve its conditions instead of assuming them to simplify the output.

\subsection{Numerical output and \code{N}}
There is no order-preserving embedding of an ordered field containing a positive infinitesimal into $\R$ that fixes the rationals. If $0<\epsilon<1/n$ for all positive integers $n$, its image would be a positive real smaller than every $1/n$, impossible by the Archimedean property. Arbitrary precision does not change this argument.

Three different numerical requests are nevertheless useful. \emph{Coefficient approximation} approximates ordinary coefficients while retaining exact monomials. \emph{Standard part} extracts the ordinary constant coefficient of a finite value, discarding infinitesimals by definition. \emph{Scale specialization} replaces finitely many monomials by selected real parameters for a numerical experiment. The last operation is not a faithful model of all surreal inequalities: no fixed positive real $u$ stays below every positive power of a fixed $0<t<1$.

A generic \code{N[value]} rule should therefore not pretend to return an ordinary complex approximation to an arbitrary surreal number. It should either expose a documented coefficient-approximation mode or require an explicit specialization. Marking every custom surreal head \code{NumericQ} can cause numerical solvers to request a conventional machine or arbitrary-precision number that the object cannot supply. The algebraic interface and the numerical interface should be designed separately.

\subsection{Serialization, caching, and complexity}
Serialize semantic constructors and parent identifiers, not just pretty output. Coefficient programs and function definitions may contain executable Wolfram code; loading untrusted serialized programs is a code-execution risk. A production format should use a validated expression grammar and explicitly trusted extension modules rather than automatically evaluating arbitrary input strings.

Caches should be keyed by normalized parent, exponent, branch, derivation, and extension policy. A coefficient computed under one phase convention cannot be reused under another. Memoized normal forms should distinguish an exact result from a precision-limited result and should never silently upgrade a guessed leading term to a certificate.

Sparse multiplication costs at most the pairwise product of input support sizes before collection; rational simplification can have severe intermediate expression growth. A generic finite-grid coefficient routine can face a large Diophantine box, while finite-algebra computations scale with the quotient dimension and coefficient complexity. Recursive exponent comparison may dominate an otherwise small calculation. These considerations favor a compact expression DAG, delayed expansion, and specialized exact subdomains over eager conversion to full normal form.

\section{Worked examples and the executable core}\label{sec:examples}
\subsection{Using the rational-monomial package}
The file \code{code/HahnRational.wl} implements \cref{thm:rationalcore} for Gaussian rational coefficients and integer exponent vectors. With the variable list \code{\{u,t\}}, the interpretation is
\[
 v(u)=(1,0),\quad v(t)=(0,1),\quad
 u=t^\omega=\omega^{-\omega},\quad t=\omega^{-1}.
\]
The leftmost coordinate is most significant. Larger valuation means smaller magnitude. Fresh, unassigned symbols should be used as monomial placeholders.

\par\noindent\begin{minipage}{\textwidth}
\begin{lstlisting}
Get["code/HahnRational.wl"];
Clear[u, t];
a = HCreate[{u, t}, 1/(1-t) + u];
b = HCreate[{u, t}, 1/(1-t)];
HValuation[HAdd[a, HScale[b, -1]]]
(* {1, 0} *)
HCompare[HCreate[{u, t}, u], HCreate[{u, t}, t^100]]
(* -1 *)
\end{lstlisting}
\end{minipage}\par
The first operation cancels an entire infinite geometric block exactly and reveals a term at a new valuation rank. A coefficient stream that had no way to recognize the equal geometric tails could fail to reach it. The second comparison is exact; it does not depend on choosing a large real value for $\omega$.

The package stores a normalized polynomial numerator and denominator, not an expanded Hahn series. \code{HValuation} computes their least exponent vectors and subtracts. \code{HSign} first verifies that an input equals its coefficientwise conjugate. \code{HCompare} rejects nonreal inputs even when their difference happens to be real. Different ordered variable lists are different parents and are not silently mixed.

The scalar domain is intentionally small. \code{HCreate[\{u,t\},Pi+t]}, a floating-point coefficient, \code{Exp[t]}, and \code{Sqrt[t]} are rejected by this backend. The first belongs in a larger coefficient parent, the last in an exponent or algebraic extension, and the exponential in a functional or lazy-series extension. Rejection means that the core does not implement that representation, not that the corresponding exact surreal value cannot be defined.

\subsection{Conjugation, squared modulus, and standard part}
For $z=t+iu$,
\[
 \overline z=t-iu,\qquad z\overline z=t^2+u^2.
\]
For a finite value with nonzero leading exponent zero, standard part is its initial coefficient. Thus
\[
 \st\!\left(\frac{2+3i+t}{1-u}\right)=2+3i,
 \qquad \st\!\left(\frac{u}{t^{20}}\right)=0.
\]
The quotient $t/u$ is infinite and has no finite standard part in this convention. The package returns a \code{Failure} for that request rather than a floating-point infinity.

\par\noindent\begin{minipage}{\textwidth}
\begin{lstlisting}
z = HCreate[{u, t}, t + I u];
HExpression[HNormSquared[z]]
(* t^2 + u^2 *)
HStandardPart[HCreate[{u, t}, (2+3 I+t)/(1-u)]]
(* 2 + 3 I *)
HStandardPart[HCreate[{u, t}, t/u]]
(* Failure["InfiniteValue", ...] *)
\end{lstlisting}
\end{minipage}\par
The internal \code{HR} head is not a public constructor. Directly forging its arguments bypasses validation and lies outside the prototype contract. Arbitrary malformed inputs are not handled as a production-grade general symbolic language; the documented operations expect successfully constructed values.

\subsection{An exact value that cannot be reached by a requested finite prefix}
The geometric remainder in \eqref{eq:geometricremainder} is retained exactly by
\par\noindent\begin{minipage}{\textwidth}
\begin{lstlisting}
HCreate[{u, t}, 1/(1-t) - Sum[t^j, {j, 0, 9}]]
\end{lstlisting}
\end{minipage}\par
Its rational value is $t^{10}/(1-t)$ and its valuation is $(0,10)$. The requested error scale $u$ has valuation $(1,0)$, strictly larger than $(0,N)$ for every ordinary integer $N$. A correct finite-prefix truncation routine must say that the requested cut needs an infinite block, while a compressed truncation routine can return the rational block exactly.

This example supplies a useful regression test for any future precision module: increasing an ordinary integer precision parameter must never claim to certify accuracy beyond $u$. The valuation-vector comparison, not a numerical norm chosen for convenience, decides whether a claimed bound is valid.

\subsection{Local analytic evaluation with ordinary coefficient tools}
The infinitesimal Taylor value of $\exp(t+t^2)$ begins
\begin{equation}\label{eq:expseriesexample}
 1+t+\frac32t^2+\frac76t^3+\frac{25}{24}t^4+O_v(5).
\end{equation}
For the rank-one parameter $t$, ordinary \code{Series} can compute these coefficients, and the strong positive-support interpretation justifies their use at $t=\omega^{-1}$. The displayed truncation is not the entire value; an exact representation would retain \code{ExpInf[t+t\^{}2]} or its certified coefficient rule. The logarithm of the exact exponential is $t+t^2$ in the infinitesimal exponential--logarithm group. The corresponding finite coefficient calculation is checked in the accompanying Python script.

For the radius-free example \eqref{eq:radiusfree}, substituting $a=t^2$ gives
\begin{equation}\label{eq:radiusfreesubstitution}
 H(t^2)=\sum_{n\ge1}\frac{t^n}{1-nt^2}
 =t+t^2+2t^3+3t^4+5t^5+9t^6+16t^7+31t^8+\cdots.
\end{equation}
At exponent $m$, the coefficient is the finite integer sum
\begin{equation}\label{eq:radiusfreecoeff}
 [t^m]H(t^2)=
 \sum_{\substack{n\ge1,\ k\ge0\\n+2k=m}}n^k.
\end{equation}
The rational coefficient germs in $z$ have no common ordinary radius, but this particular coefficient algorithm terminates with an explicit finite dependency set. CT identifies the same value as a microscopic Cauchy integral at an appropriate radius \cite{CT}. The minimal rational-field package does not represent the outer infinite sum; \eqref{eq:radiusfreecoeff} describes a straightforward next lazy-series module.

\subsection{A finite algebra that retains four singularly perturbed points}
Consider the polynomial system from AG and CT,
\begin{equation}\label{eq:fourexample}
 F_1=x^2-Ay,\qquad F_2=y^2-Bx,
 \qquad A,B\ne0\text{ infinitesimal}.
\end{equation}
With basis $(1,x,y,xy)$ and $q=AB$, its coordinate multiplication matrices are
\begin{equation}\label{eq:fourmatrices}
 M_x=\begin{pmatrix}0&0&0&0\\1&0&0&q\\0&A&0&0\\0&0&1&0\end{pmatrix},
 \qquad
 M_y=\begin{pmatrix}0&0&0&0\\0&0&B&0\\1&0&0&q\\0&1&0&0\end{pmatrix}.
\end{equation}
They commute and satisfy $M_x^2=AM_y$ and $M_y^2=BM_x$. Their characteristic polynomials are
\[
 Q_x(T)=T(T^3-A^2B),\qquad
 Q_y(T)=T(T^3-AB^2).
\]
All these identities are finite polynomial calculations. Taking $A=t$ and $B=u=t^\omega$ preserves the higher-rank distinction without changing the matrix algorithms.

The four roots are the origin and
\[
 (\rho\zeta^j,\rho^2\zeta^{2j}/A),\qquad
 j=0,1,2,\qquad \rho^3=A^2B,
\]
where $\zeta$ is a primitive cube root of unity. Their nonzero coordinate valuations are $(2v(A)+v(B))/3$ and $(v(A)+2v(B))/3$. These roots require an algebraic/exponent extension even though the multiplication matrices do not. Keeping the finite algebra avoids forcing that extension before computing traces or total residues \cite{AG,CT}.

The Jacobian is $J_F=4xy-AB$. The simple residue weights for the constant numerator are
\[
 -\frac1{AB},\quad\frac1{3AB},\quad\frac1{3AB},\quad\frac1{3AB}.
\]
Their sum is exactly zero although each individual weight is infinite. Meanwhile $\Res_F(xy)=1$ and $\Res_F(J_F)=4$. A truncation or floating-point specialization that loses cancellation can therefore fail on a quantity with an especially simple exact answer.

For a monomial numerator, the residue is
\begin{equation}\label{eq:monomialresidue}
 \Res_F(x^p y^r)=
 \begin{cases}
 A^{(2p+r-3)/3}B^{(p+2r-3)/3},&
 (2p+r-3)/3,(p+2r-3)/3\in\N,\\
 0,&\text{otherwise}.
 \end{cases}
\end{equation}
The independent checks compare the trace identity and the separating-denominator computation with \eqref{eq:monomialresidue} for all $0\le p,r\le10$.

\subsection{Why the separating determinant matters}
For \eqref{eq:fourexample}, set
\[
 \mathsf C=\begin{pmatrix}x^2+Ay&A^2\\B^2&y^2+Bx\end{pmatrix}.
\]
Then
\[
 \mathsf C\binom{F_1}{F_2}
 =\binom{x^4-A^2Bx}{y^4-AB^2y},
\]
and
\[
 \det\mathsf C=x^2y^2+Bx^3+Ay^3+ABxy-A^2B^2.
\]
CT's formula integrates $H\det\mathsf C/(Q_x(x)Q_y(y))$ over a separating torus. The separated system has sixteen simple grid points, but the original system has four. At a grid point not solving the original system, $\mathsf C F=0$ with $F\ne0$, so $\det\mathsf C=0$ and the extraneous simple residue vanishes. Dropping the determinant would compute the wrong invariant \cite{CT}.

At the algebraic level, expansion of the reciprocals gives
\[
 Q_x(x)^{-1}=x^{-4}\sum_{k\ge0}(A^2B)^kx^{-3k},\qquad
 Q_y(y)^{-1}=y^{-4}\sum_{\ell\ge0}(AB^2)^\ell y^{-3\ell}.
\]
On an admissible scaled torus these expansions are strongly supported. Extracting $x^{-1}y^{-1}$ after multiplying by $H\det\mathsf C$ reproduces \eqref{eq:monomialresidue}. The code checks this finite coefficient calculation; it does not numerically integrate a fine-topological torus.

\subsection{Three circles that a single real precision can confuse}
In $\Gamma=\Q\omega+\Q$, take
\[
 R(z)=\frac1{z-t}+\frac1{z-u},\qquad u=t^\omega.
\]
A positively oriented microscopic circle $|z|=t^s$ has normalized period
\[
 \frac1{2\pi i}\HInt_{|z|=t^s}R(z)\,dz
 =\begin{cases}2,&s=1/2,\\1,&s=2,\\0,&s=2\omega.\end{cases}
\]
A pole $t^a$ is inside exactly when $a>s$; both residues equal one. This reduces the computation to two exact exponent comparisons. The example is from CT and is independently checked at the valuation-vector level in the accompanying script \cite{CT}. The rational-monomial core already supplies the essential order operations, while a complete contour wrapper would still need its chain and pole-avoidance certificates.

\section{Implementation priorities and an operation matrix}\label{sec:roadmap}
\subsection{A staged development plan}
A useful first release should extend the tested rational-monomial core rather than start with arbitrary surreal cuts. Its immediate targets are stronger validation, efficient canonical fraction data, algebraic coefficient fields, rational exponent lattices, explicit parent embeddings, and a stable expression/serialization format. Those tasks enlarge a domain where exact equality and real order remain reliable.

The next stage should add selected algebraic roots and rank-one or nested Puiseux computations, together with certified finite precision. Root descriptors must carry selection data, and refinement must preserve branch identity. Higher-rank precision should permit exact blocks and explicit unresolved cuts, not just an integer truncation order.

A third stage can add positive-grid lazy series, infinitesimal Taylor lifting, local implicit functions, and analytic-germ interfaces. Its first supported coefficient functions can be rational and algebraic germs, followed by selected ordinary-point differential descriptions. Implement \cref{prop:grid}'s dependency guarantee before offering unrestricted-looking sum and composition syntax.

A fourth stage can build finite quotient algebras, polynomial complete-intersection deformation, residue pairings, and rational contour periods. These are rich operations with compact certificates. The arbitrary radius-free analytic-germ generalization should follow only when an effective ordinary-germ division backend exists for the chosen language.

Finally, connect a genuine transseries backend for real exponential--logarithmic and differential-algebraic extensions. Surcomplex phase policies, sectorial summation, and global chart compatibility are distinct projects. They should not delay a usable exact algebraic implementation, nor be marked implemented merely because the parser accepts their names.

\subsection{What is exact, decidable, or merely conditional?}
In the following table, ``effective'' always includes the stated coefficient and exponent operations. ``Conditional'' means that support, branch, coefficient-algorithm, or category hypotheses must be supplied and checked; it is not a claim of a general terminating algorithm.
\begin{center}\small
\begin{longtable}{@{}L{0.27\textwidth}L{0.31\textwidth}L{0.33\textwidth}@{}}
\toprule Operation & Effective rational/algebraic layer & Broader series/function layer\\\midrule\endhead
Addition, multiplication & Exact finite arithmetic; decidable equality & Strong convolution; coefficient access needs effective dependencies.\\[0.4em]
Inverse & Exact for a proved nonzero rational value & Normalize a proved leading term; strong geometric inversion or selected algebraic inverse.\\[0.4em]
Equality and real order & Total for the rational core; effective selected real-algebraic extensions & Partial in general; total in supported structured classes, including appropriate transseries zero-test domains.\\[0.4em]
Valuation and leading coefficient & Exact finite polynomial calculation & Cannot be total for arbitrary coefficient programs.\\[0.4em]
Coefficient at one exponent & Finite expansions or specialized rational algorithms & Effective positive grids and nested supported descriptions; arbitrary supports require extra data.\\[0.4em]
All terms below a cut & May return an exact compressed rational object & A finite list may not exist, even for a valid Hahn series.\\[0.4em]
Real/imaginary parts, conjugation & Exact in the complex rational core & Coefficientwise when supported.\\[0.4em]
Modulus and algebraic roots & Selected real/algebraic extension may be required & Exact root descriptor and branch selection; finite normal form not guaranteed.\\[0.4em]
Standard part & Exact on finite inputs with known leading data & May be available from a coefficient query plus a finiteness certificate.\\[0.4em]
Taylor evaluation & Polynomial and rational substitution & Exact lazy lift at supported centers and infinitesimal displacements.\\[0.4em]
Exponential and logarithm & Usually leave the rational parent & Local infinitesimal series; real transseries extension; surcomplex phase and branch policies.\\[0.4em]
Differentiation & Exact formal-variable rules & Must specify coefficient derivative, parameter derivative, or surreal derivation.\\[0.4em]
Integration and residues & Rational exactness/residue algorithms & Chosen coherent or derivation-based theory; no universal fine-local fundamental theorem.\\[0.4em]
Root counts and finite schemes & Polynomial algebra, matrices, selected roots & Manuscript deformation methods plus effective ordinary division and support.\\[0.4em]
Numerical output & Coefficient approximation or explicit specialization & No ordinary real approximation preserving all infinitesimal scales.\\
\bottomrule
\end{longtable}\end{center}

\subsection{Acceptance tests for future modules}
A future algebraic-root module should verify defining equations, branch separation, parent embeddings, and multiplicities, including repeated-root cases. A support module should reject positive but non-well-ordered supports and well-ordered unions with infinitely many contributions to one exponent. A precision module should preserve the counterexample in \eqref{eq:geometricremainder}. A function module should distinguish $\exp(t)$ from an invalid strong expansion of $\exp(1/t)$ and should preserve branch conditions in logarithmic rewrites.

A finite-algebra module should pass the four-point example, the overdetermined length-drop counterexample, and the sharp Hensel threshold example. A contour module should detect poles on the contour, use the transformation determinant, and distinguish ordinary-parameter Hahn chains from fine-continuous paths. A global surcomplex exponential module should make the two values in \eqref{eq:phasecontrast} explicit under the two named policies, never cache them as if they belonged to one function.

These tests do not replace proofs of the general algorithms, but they target errors that ordinary complex arithmetic tests would miss. A proof-oriented interface could export finite polynomial identities, support derivations, valuation inequalities, and root-selection certificates to a checker. Such a checker would verify specific claims without needing to decide every question about every surreal object.

\section{Conclusions}\label{sec:conclusions}
A substantial exact symbolic surreal CAS is possible without a universal representation of all surreals. The most reliable core consists of explicit effective coefficient fields, ordered exponent groups, finite Hahn sums, and rational monomial fractions. Selected algebraic roots add exact polynomial solving and modulus. Certified lazy supports and infinitesimal Taylor lifts add a large local function language. Structured transseries supply a principled route toward infinite-argument exponential--logarithmic and differential-algebraic computation.

The main limitations have different sources and should not be conflated. Cardinality prevents a finite language from naming every real or surreal. Undecidability prevents a total zero-test for arbitrary coefficient programs. Higher-rank support geometry prevents some finite-prefix precision requests from having finite answers. Analytic domains and branch policies prevent unrestricted substitutions and simplifications. The supplied manuscripts further show that contour theorems and global transcendental functions depend on the chosen analytic category.

The constructive lesson is to attach the necessary information to the object. A parent identifies arithmetic and order. A support presentation identifies coefficient dependencies. A root descriptor identifies a branch. A derivative identifies its constants. A contour identifies a chain category and pole-avoidance conditions. A global surcomplex function identifies its extension policy. With those interfaces, the strong formulas of the supplied manuscripts become either genuine algorithms on effective inputs or clearly labeled conditional constructions. Without them, elegant symbolic notation can conceal undefined operations or nonterminating searches.

The accompanying prototype demonstrates the central algebraic point: exact multi-rank arithmetic does not require truncating infinitesimals or replacing infinity by a numerical limit. The more ambitious analytic and transseries layers are specified here as extensions, with their mathematical and computational obligations made explicit.

\appendix
\section{Reproducibility and what the checks establish}\label{app:repro}
The archive contains the article source and PDF, the package and example files in the \code{code/} directory, and two recorded validation reports. The executable files are \path{HahnRational.wl}, \path{test_HahnRational.wl}, \path{examples.wl}, and \path{verify_examples.py}. The package's executable definitions and test body were loaded through \code{Get[StringToStream[...]]} in the connected Wolfram kernel, preserving sequential package-context parsing. The recorded clean result was
\par\noindent\begin{minipage}{\textwidth}
\begin{lstlisting}
<|"Kernel" -> "15.0.1 for Linux x86 (64-bit) (July 2, 2026)",
  "Tests" -> 88, "Passed" -> 88, "Failures" -> {}|>
\end{lstlisting}
\end{minipage}\par
The test file includes deterministic examples and seeded random polynomial inputs. It tests exact rational field identities, higher-rank valuation and ordering, conjugation, standard part, and rejection of unsupported operations. The seed is \code{20260921}. To reproduce the tests with a local Wolfram installation, run
\par\noindent\begin{minipage}{\textwidth}
\begin{lstlisting}
wolframscript -file code/test_HahnRational.wl
\end{lstlisting}
\end{minipage}\par
The kernel version is a property of the executed test environment, not an assertion that every earlier or later Wolfram release has been tested.

The independent Python script, executed with Python 3.13.5 and SymPy 1.14.0, passed \textbf{323 exact checks}. These include the multiplication relations, characteristic polynomials, residue-pairing matrix, separating-denominator identities, trace and torus residue calculations on a $11\times11$ monomial grid, radius-free coefficients through exponent sixteen, geometric remainders, infinitesimal exponential/logarithm coefficients, the factorial-series differential equation, and the three valuation-based contour counts. Run
\par\noindent\begin{minipage}{\textwidth}
\begin{lstlisting}
python code/verify_examples.py
\end{lstlisting}
\end{minipage}\par
The two suites check different things; their counts should not be interpreted as a formal proof of 411 general theorems. The Wolfram package is not an implementation of arbitrary Hahn fields, algebraic closure, general lazy series, AG's arbitrary-germ division, the transseries zero-test, or CT's general integration theory. The Python script verifies finite examples independently; it is not a replacement implementation of those theories.

\section{Source map and boundaries of reliance}\label{app:sources}
\begin{center}\small
\begin{longtable}{@{}L{0.25\textwidth}L{0.66\textwidth}@{}}
\toprule Input & Use in this article\\\midrule\endhead
SA, supplied manuscript & Fine-local versus coherent functions; strong local evaluation; phase-dependent exponential extensions; logarithm category warning; all-scale rigidity. These are attributed manuscript results.\\[0.4em]
AG, supplied manuscript & Radius-free germ algebra; evaluation and translation; division and normal forms; finite-algebra deformation, residues, traces, stability, and effectivity warning. General theorems are not independently machine-verified here.\\[0.4em]
CT, supplied manuscript & Coherent chain integration; microscopic rescaling; separating-torus formula with determinant; algebraic rational periods and worked contour examples. The stated category hypotheses are retained.\\[0.4em]
Published surreal/transseries literature & Foundational constructions, differential structures, structured exact function languages, and relevant current zero-testing developments. No claim that one cited framework equals every other.\\[0.4em]
Official CAS documentation & Contracts of specific Wolfram and Sage facilities. These motivate reuse, not a claim that the facilities implement the proposed surreal semantics already.\\[0.4em]
Arguments and code in this article & Effective rational-monomial core, finite-grid dependency proof, representation-sensitive halting obstruction, precision derivations, proposed interfaces, and executed finite tests. No exhaustive novelty or priority claim is intended.\\
\bottomrule
\end{longtable}\end{center}
The mathematical review and documentation checks were conducted for this article on September 21, 2026. In particular, the ISSAC 2026 zero-test is included as a positive algorithmic result rather than overlooked because older general discussions emphasize undecidability. A complete implementation would need a separate engineering and proof review of every imported algorithm.

\begin{thebibliography}{99}
\addcontentsline{toc}{section}{References}
\bibitem{SA}
\emph{Surcomplex Analysis: Infinitesimal Calculus, Hahn-Coherent Holomorphy, and Contour Theory}.
Supplied research manuscript, September 21, 2026.
Source file: \code{surcomplex\_analysis(3).tex}. Especially Sections 2--6, 11--13.

\bibitem{AG}
\emph{Infinitesimal Analytic Geometry over the Surcomplex Numbers: Radius-Free Hahn Germs, a Nullstellensatz, and Conservation of Singular Zeros}.
Supplied research manuscript, September 21, 2026.
Source file: \code{surcomplex\_analytic\_geometry.tex}. Especially Sections 2--4, 6--12 and Appendix B.

\bibitem{CT}
\emph{Jordan Separation, Cauchy Theory, and Stokes Formulae on the Surcomplex Plane: A Category-Sensitive Contour Calculus}.
Supplied research continuation, September 21, 2026.
Source file: \code{surcomplex\_contours(2).tex}. Especially Sections 2, 4--10 and Appendices A--C.

\bibitem{Gonshor}
Harry Gonshor, \emph{An Introduction to the Theory of Surreal Numbers}.
London Mathematical Society Lecture Note Series 110, Cambridge University Press, 1986.
\href{https://doi.org/10.1017/CBO9780511629143}{doi:10.1017/CBO9780511629143}.

\bibitem{Poonen}
Bjorn Poonen, ``Maximally complete fields,''
\emph{L'Enseignement Math\'ematique} (2) \textbf{39} (1993), 87--106.
\href{https://math.mit.edu/~poonen/papers/amsval.pdf}{Author's full text}.
Corollaries 4 and 7 concern algebraic closedness and the characteristic-zero Puiseux extension.

\bibitem{vdH}
Joris van der Hoeven, \emph{Transseries and Real Differential Algebra}.
Lecture Notes in Mathematics 1888, Springer, 2006.
\href{https://doi.org/10.1007/3-540-35590-1}{doi:10.1007/3-540-35590-1};
\href{https://www.texmacs.org/joris/ln/ln.html}{Author's electronic text}.

\bibitem{CFH}
Shaoshi Chen, Hanqian Fang, and Joris van der Hoeven,
``A Zero-Test for D-Algebraic Transseries,''
in \emph{Proceedings of ISSAC 2026}, ACM, 2026, pp.~105--113.
Published July 12, 2026.
\href{https://doi.org/10.1145/3815436.3815443}{doi:10.1145/3815436.3815443};
\href{https://arxiv.org/abs/2602.01188}{arXiv:2602.01188};
\href{https://www.texmacs.org/joris/zttr/zttr.html}{Author's full text}.
Especially Sections 2.4--2.5 and 7, including Theorem 18.

\bibitem{BMder}
Alessandro Berarducci and Vincenzo Mantova,
``Surreal numbers, derivations and transseries,''
\emph{Journal of the European Mathematical Society} \textbf{20} (2018), 339--390.
\href{https://doi.org/10.4171/JEMS/769}{doi:10.4171/JEMS/769};
\href{https://arxiv.org/abs/1503.00315}{arXiv:1503.00315}.

\bibitem{BMfun}
Alessandro Berarducci and Vincenzo Mantova,
``Transseries as germs of surreal functions,''
\emph{Transactions of the American Mathematical Society} \textbf{371} (2019), 3549--3592.
\href{https://arxiv.org/abs/1703.01995}{arXiv:1703.01995}.

\bibitem{BG}
Olivier Bournez and Quentin Guilmant,
``Surreal fields stable under exponential, logarithmic, derivative and anti-derivative functions,''
preprint, 2022.
\href{https://arxiv.org/abs/2211.08396}{arXiv:2211.08396}.

\bibitem{CE}
Ovidiu Costin and Philip Ehrlich,
``Integration on the surreals,'' \emph{Advances in Mathematics}
\textbf{452} (2024), article 109823.
\href{https://doi.org/10.1016/j.aim.2024.109823}{doi:10.1016/j.aim.2024.109823};
\href{https://arxiv.org/abs/2208.14331}{arXiv:2208.14331}.

\bibitem{CDS}
Eduardo Cattani, Alicia Dickenstein, and Bernd Sturmfels,
``Computing multidimensional residues,'' in
L.~Gonz\'alez-Vega and T.~Recio (eds.),
\emph{Algorithms in Algebraic Geometry and Applications},
Progress in Mathematics 143, Birkh\"auser, 1996, pp.~135--164.
\href{https://doi.org/10.1007/978-3-0348-9104-2_8}{doi:10.1007/978-3-0348-9104-2\_8};
\href{https://arxiv.org/abs/alg-geom/9404011}{arXiv:alg-geom/9404011}.

\bibitem{WLRoot}
Wolfram Research, \emph{Root}, Wolfram Language documentation.
\url{https://reference.wolfram.com/language/ref/Root.html}.
Accessed September 21, 2026.

\bibitem{WLAlgebraic}
Wolfram Research, \emph{AlgebraicNumber} and \emph{Algebraic Numbers}, Wolfram Language documentation.
\url{https://reference.wolfram.com/language/ref/AlgebraicNumber.html};
\url{https://reference.wolfram.com/language/guide/AlgebraicNumbers.html}.
Accessed September 21, 2026.

\bibitem{WLSeries}
Wolfram Research, \emph{SeriesData}, Wolfram Language documentation.
\url{https://reference.wolfram.com/language/ref/SeriesData.html}.
Accessed September 21, 2026.

\bibitem{WLAsymptotic}
Wolfram Research, \emph{Asymptotic}, Wolfram Language documentation.
\url{https://reference.wolfram.com/language/ref/Asymptotic.html}.
Accessed September 21, 2026.

\bibitem{WLInfinity}
Wolfram Research, \emph{Infinity}, Wolfram Language documentation.
\url{https://reference.wolfram.com/language/ref/Infinity.html}.
Accessed September 21, 2026.

\bibitem{WLInactive}
Wolfram Research, \emph{Inactive}, Wolfram Language documentation.
\url{https://reference.wolfram.com/language/ref/Inactive.html}.
Accessed September 21, 2026.

\bibitem{WLUpValues}
Wolfram Research, \emph{UpValues}, Wolfram Language documentation.
\url{https://reference.wolfram.com/language/ref/UpValues.html}.
Accessed September 21, 2026.

\bibitem{SageLazy}
The Sage Developers, \emph{Lazy Series}, SageMath reference manual.
\href{https://doc.sagemath.org/html/en/reference/power_series/sage/rings/lazy_series.html}{Official lazy-series reference documentation}.
Accessed September 21, 2026.

\bibitem{SageAsymptotic}
The Sage Developers, \emph{Asymptotic Ring}, SageMath reference manual.
\href{https://doc.sagemath.org/html/en/reference/asymptotic/sage/rings/asymptotic/asymptotic_ring.html}{Official asymptotic-ring reference documentation}.
Accessed September 21, 2026.
\end{thebibliography}
\end{document}
